\chapter{ES Optimization Implementation} \label{app:qp_formulations}

Building upon the continuous equivalence derived in Appendix \ref{app:es_derivation} and following the optimization framework in the original paper (Ulyrch, Lucescu 2026), this section details the precise matrix formulations required to solve the regularized multi-currency Expected Shortfall models. By linearizing the $L_1$ penalty and discretizing the ES computations, we map both the separate overlay and joint optimization frameworks into standard Quadratic Programs of the form
\begin{equation} \label{eq:standard_qp}
    \min_{\tilde{z}} \left( \tilde{c}^T \tilde{z} + \frac{1}{2} \tilde{z}^T \tilde{Q} \tilde{z} \right) \quad \text{subject to} \quad \tilde{D} \tilde{z} = \tilde{E}, \quad \tilde{A} \tilde{z} \le b, \quad \tilde{z} \ge 0,
\end{equation}
where $\tilde{z}$ is an augmented decision vector containing the portfolio weights and all necessary auxiliary variables.

\section{Sample-Based Expected Shortfall}

In practice, the conditional expectation from the ES objective is approximated using $q$ historical joint return observations. Let $R_{\text{data}} \in \mathbb{R}^{q \times N}$ be the matrix of historical asset returns, where each row $R_k^T$ represents the return vector at time $k$. The sample-based objective becomes:
\begin{equation*}
    \min_{x_t, \alpha} \left( - \mu^T x_t + \gamma \alpha + \frac{\gamma}{q(1-\beta)} \sum_{k=1}^q (-R_k^T x_t - \alpha)^+ \right).
\end{equation*}

To linearize the maximum operator $(\cdot)^+$, Rockafellar and Uryasev (2000) introduce a $q$-dimensional vector of auxiliary variables $u \in \mathbb{R}^q$, where $u_k = (-R_k^T x_t - \alpha)^+$. This implies two conditions for each sample $k$:
\begin{align*}
    u_k &\ge 0 \\
    u_k &\ge -R_k^T x_t - \alpha.
\end{align*}
Rearranging the second condition to isolate the variables yields a standard linear inequality:
\begin{equation} \label{eq:es_inequality}
    - R_k^T x_t - \alpha - u_k \le 0 \quad \forall \, k \in \{1, \dots, q\}.
\end{equation}
In block matrix notation across all $q$ samples, this becomes:
\begin{equation*}
    -R_{\text{data}} x_t - \mathbf{1}_q \alpha - I_q u \le \mathbf{0}_q,
\end{equation*}
where $\mathbf{1}_q$ is a $q$-dimensional vector of ones, $I_q$ is the $q \times q$ identity matrix, and $\mathbf{0}_q$ is a $q$-dimensional vector of zeros. This constructs the first three block columns of the general inequality constraint matrix $\tilde{A}$.

\section{Linearization of the $L_1$ Penalty and Leverage}

The $L_1$ norm $\|x_t\|_1 = \sum_{i=1}^N |x_{i,t}|$ introduces non-differentiability into the objective function. To resolve this, we introduce non-negative auxiliary variables representing the positive and negative components of each asset weight: $\Delta^+_t, \Delta^-_t \in \mathbb{R}^N$ such that $\Delta^+_t \ge 0$ and $\Delta^-_t \ge 0$. 

The weights are decomposed as:
\begin{equation} \label{eq:weight_decomp}
    x_t = \Delta^+_t - \Delta^-_t \implies I_N x_t + I_N \Delta^-_t - I_N \Delta^+_t = \mathbf{0}_N.
\end{equation}
Because the optimization minimizes the penalties, at most one of $\Delta^+_{i,t}$ or $\Delta^-_{i,t}$ will be non-zero for any asset $i$. Therefore, the absolute value is exactly determined by their sum: $|x_{i,t}| = \Delta^+_{i,t} + \Delta^-_{i,t}$.

This allows the $L_1$ penalty term to be formulated linearly in the objective function as $\lambda_1 \mathbf{1}_N^T \Delta^+_t + \lambda_1 \mathbf{1}_N^T \Delta^-_t$. Furthermore, the portfolio leverage constraint $\|x_t\|_1 \le \ell$ transforms into a linear inequality:
\begin{equation} \label{eq:leverage_inequality}
    \mathbf{1}_N^T \Delta^+_t + \mathbf{1}_N^T \Delta^-_t \le \ell.
\end{equation}

\section{Quadratic Translation of the $L_2$ Penalty}

Unlike the expected shortfall and $L_1$ terms which map to the linear objective vector $\tilde{c}$, the $L_2$ penalty integrates naturally into the quadratic matrix $\tilde{Q}$. The penalty is defined as:
\begin{equation*}
    \lambda_2 \|x_t\|_2^2 = \lambda_2 x_t^T x_t = \frac{1}{2} x_t^T (2 \lambda_2 I_N) x_t.
\end{equation*}
Thus, the block of $\tilde{Q}$ corresponding to the weights $x_t$ is populated with the diagonal matrix $2 \lambda_2 I_N$, while all auxiliary variables receive zeros.

\section{Separate Optimization (Currency Overlay)}

In the separate approach, asset weights $x_t \in \mathbb{R}^N$ and currency exposures $\psi_t \in \mathbb{R}^M$ are optimized sequentially.

\subsection{Step 1: Asset-Level Optimization}
The target vector relies on fully hedged asset returns $R^{\text{fh}} \in \mathbb{R}^{q \times N}$. We define the augmented vector of dimension $(3N + q + 1)$:
\begin{equation*}
    \tilde{x}_t = \begin{pmatrix} x_t^T, & \alpha, & u^T, & (\Delta_t^-)^T, & (\Delta_t^+)^T \end{pmatrix}^T.
\end{equation*}

The objective vector $\tilde{c}$ and quadratic matrix $\tilde{Q}$ are constructed as:
\begin{equation*}
    \tilde{c} = \begin{pmatrix} -\mu_{\text{fh}} \\ \gamma \\ \frac{\gamma}{q(1-\beta)} \mathbf{1}_q \\ \lambda_1^a \mathbf{1}_N \\ \lambda_1^a \mathbf{1}_N \end{pmatrix}, \quad 
    \tilde{Q} = \begin{pmatrix} 2\lambda_2^a I_N & O \\ O & O \end{pmatrix},
\end{equation*}
where $O$ represents zero matrices of appropriate dimensions. The equality constraints enforce the budget condition $\mathbf{1}_N^T x_t = 1$ and the variable decomposition from Equation \ref{eq:weight_decomp}:
\begin{equation*}
    \tilde{D} = \begin{pmatrix} 
    \mathbf{1}_N^T & 0 & \mathbf{0}_q^T & \mathbf{0}_N^T & \mathbf{0}_N^T \\ 
    I_N & \mathbf{0}_N & O_{N \times q} & I_N & -I_N 
    \end{pmatrix}, \quad 
    \tilde{E} = \begin{pmatrix} 1 \\ \mathbf{0}_N \end{pmatrix}.
\end{equation*}

The inequality matrix $\tilde{A}$ and bound vector $b$ encode the shortfall losses from Equation \ref{eq:es_inequality} and the leverage limit from Equation \ref{eq:leverage_inequality}:
\begin{equation*}
    \tilde{A} = \begin{pmatrix} 
    -R^{\text{fh}} & -\mathbf{1}_q & -I_q & O_{q \times N} & O_{q \times N} \\
    \mathbf{0}_N^T & 0 & \mathbf{0}_q^T & \mathbf{1}_N^T & \mathbf{1}_N^T 
    \end{pmatrix}, \quad 
    b = \begin{pmatrix} \mathbf{0}_q \\ \ell \end{pmatrix}.
\end{equation*}

\subsection{Step 2: Currency-Level Optimization}
Conditional on the optimal asset weights $x_t^*$, the currency exposures $\psi_t$ are optimized. Using historical excess currency returns $R^c \in \mathbb{R}^{q \times M}$, the augmented vector becomes $\tilde{\psi}_t \in \mathbb{R}^{(3M + q + 1)}$.

Because currency overlays act as zero-investment portfolios, there is no budget constraint. Additionally, the pre-fixed asset portfolio returns ($R^{\text{fh}} x_t^*$) must be subtracted from the shortfall inequality constraint boundaries. Thus, the formulations shift to:
\begin{equation*}
    \tilde{D}_{\psi} = \begin{pmatrix} 
    I_M & \mathbf{0}_M & O_{M \times q} & I_M & -I_M 
    \end{pmatrix}, \quad \tilde{E}_{\psi} = \mathbf{0}_M,
\end{equation*}
\begin{equation*}
    \tilde{A}_{\psi} = \begin{pmatrix} 
    -R^c & -\mathbf{1}_q & -I_q & O_{q \times M} & O_{q \times M} 
    \end{pmatrix}, \quad b_{\psi} = R^{\text{fh}} x_t^*.
\end{equation*}

\section{Joint Optimization Framework}

The joint formulation merges assets and currencies into a single decision vector $\theta_t = (x_t^T, \psi_t^T)^T \in \mathbb{R}^{N+M}$, evaluated over the joint return block matrix $R^{\text{all}} = \begin{pmatrix} R^{\text{fh}} & R^c \end{pmatrix} \in \mathbb{R}^{q \times (N+M)}$.

We construct an expansive augmented vector $\tilde{\theta}_t \in \mathbb{R}^{(3N + 3M + q + 1)}$:
\begin{equation*}
    \tilde{\theta}_t = \begin{pmatrix} x_t^T, & \psi_t^T, & \alpha, & u^T, & (\Delta_{x,t}^-)^T, & (\Delta_{\psi,t}^-)^T, & (\Delta_{x,t}^+)^T, & (\Delta_{\psi,t}^+)^T \end{pmatrix}^T.
\end{equation*}

The objective vector applies separate $L_1$ penalty parameters ($\lambda_1^a$ for assets, $\lambda_1^c$ for currencies):
\begin{equation*}
    \tilde{c} = \begin{pmatrix} -\mu_{\text{fh}} \\ -\mu_c \\ \gamma \\ \frac{\gamma}{q(1-\beta)} \mathbf{1}_q \\ \lambda_1^a \mathbf{1}_N \\ \lambda_1^c \mathbf{1}_M \\ \lambda_1^a \mathbf{1}_N \\ \lambda_1^c \mathbf{1}_M \end{pmatrix}, \quad 
    \tilde{Q} = \begin{pmatrix} 
    2\lambda_2^a I_N & O & O \\ 
    O & 2\lambda_2^c I_M & O \\ 
    O & O & O 
    \end{pmatrix}.
\end{equation*}

The budget constraint applies strictly to the asset weights, while the decomposition constraints apply to both assets and currencies. 
\begin{equation*}
    \tilde{D} = \begin{pmatrix} 
    \mathbf{1}_N^T & \mathbf{0}_M^T & 0 & \mathbf{0}_q^T & \mathbf{0}_{N+M}^T & \mathbf{0}_{N+M}^T \\ 
    I_{N+M} & & \mathbf{0}_{N+M} & O_{(N+M) \times q} & I_{N+M} & -I_{N+M} 
    \end{pmatrix}, \quad 
    \tilde{E} = \begin{pmatrix} 1 \\ \mathbf{0}_{N+M} \end{pmatrix}.
\end{equation*}

Finally, the inequality matrix binds the unified shortfall distribution and enforces the leverage limit solely on the gross asset exposure:
\begin{equation*}
    \tilde{A} = \begin{pmatrix} 
    -R^{\text{all}} & -\mathbf{1}_q & -I_q & O_{q \times (2N+2M)} \\
    \mathbf{0}_{N+M}^T & 0 & \mathbf{0}_q^T & \begin{pmatrix} \mathbf{1}_N^T & \mathbf{0}_M^T \end{pmatrix} & \begin{pmatrix} \mathbf{1}_N^T & \mathbf{0}_M^T \end{pmatrix}
    \end{pmatrix}, \quad 
    b = \begin{pmatrix} \mathbf{0}_q \\ \ell \end{pmatrix}.
\end{equation*}
Solving this comprehensive quadratic program yields the optimal regularized asset and currency allocations concurrently.

\end{document}
